\section{True Diversity}

In \textbf{East and West}, Rene Guenon makes this fundamental claim:

\begin{quotationx}
So long as western people imagine that there only exists a single type of humanity, that there is only one `civilization', at different stages of development, no mutual understanding will be possible. The truth is that there are many civilizations, developing along very different lines, and that, among these, that of the modern West is strangely exceptional, as some of its characteristics show.

One should never speak absolutely of superiority or inferiority, without making it quite clear from what point of view the things to be compared are being considered, even supposing that they are comparable. There is no civilization which is superior to the others from every point of view, because man cannot be equally active at the same time in every direction, and because there are some ways of development which seem actually incompatible with one another.
\end{quotationx}

By ``stages of development'', Guenon does not mean what Westerners call ``progress''. He rejects Enlightenment notions of progress---such as Hegelianism, Marxism and similar schemes---that would force other cultures into some predefined historical pattern. Furthermore, there are groups dedicated to ``saving Western civilization'', though what they would save are the more degenerate aspects of that civilization. There is no possibility of an alliance with them.


\flrightit{Posted on 2010-08-18 by Cologero}

\begin{center}* * *\end{center}

\begin{footnotesize}
\begin{sffamily}
\texttt{kadambari on 2010-08-21 at 10:42 said:}

Evola sums up the true dignified, tolerant mind in the Doctrine of Awakening even though it is applied to Buddhism. He is completely on the mark here:

A particular characteristic of the Aryan-ness of the original Buddhist teaching is the absence of those proselytizing manias that exist, almost without exception, in direct proportion to the plebeian and anti-aristocratic character of a belief. An Aryan mind has too much respect for other people, and its sense of its own dignity is too pronounced to allow it to impose its own ideas upon others, even when it knows that its ideas are correct. Accordingly, in the original cycle of Aryan civilizations, both Eastern and Western, there is not the smallest trace of divine figures being so con¬cerned with mankind as to come near to pursuing them in order to gain their adher-ence and to ``save'' them. The so-called salvationist religions-the Erlösungsreligionen, in German-make their appearance both in Europe and Asia at a later date, together with a lessening of the preceding spiritual tension, with a fall from Olympian consciousness and, not least, with influxes of inferior ethnic and social elements. That the divinities can do little for men, that man is fundamentally the artificer of his own destiny, even of his development beyond this world-this char¬acteristic view held by original Buddhism demonstrates its difference from some later forms... 

I find the problem with ``diversity'' is that anyone can write books, does not take much be be a so-called ``expert'' on something these days, people just want to publish for a career or name, does not matter if they have anything worthwhile to say ; you have too many ``cows writing books'' (Count Keyserling, Travel Diary of a Philosopher) with the result that a great deal of confusion is generated as to what is of superior quality and what is not.

\hfill

\texttt{kadambari on 2010-08-21 at 15:52 said:}

I meant to say: does not matter if they do ``not'' have anything worthwhile to say... and ``truly dignified''... 

Also people do not want to submit to authority these days, even if it is the authority of their betters (not to mention the problem also of the deluded types who think themselves ``superior'', in the manner of Lenin's vanguard elites and cause great harm thereby).


\hfill

\end{sffamily}
\end{footnotesize}

